\begin{abstract}
We characterize minimax estimation and honest inference for a latent-cutoff treatment effect when treatment records the true cutoff side and the running variable is measured with independent Gaussian error. The benchmark has known latent support \([-1,1]\), a continuous score density bounded between \(1/4\) and \(3/4\), and binary potential outcomes whose continuous conditional means lie between the same bounds and have full-interval Hölder seminorm at most two. For each fixed Hölder exponent \(\beta\in(0,1]\), minimax expected absolute error and minimax expected length of connected intervals with uniform \(90\%\) coverage have a common rate, uniformly over sample sizes \(n\ge2\) and known noise standard deviations \(\sigma\in[0,1]\). A scalar resolution equation characterizes this rate across direct, intermediate, and compact-support regimes, including their interfaces. The direct rate is \(n^{-\beta/(2\beta+1)}\). For each fixed positive noise level, the asymptotic rate as \(n\to\infty\) is \(\{\log\log(en)/\log(en)\}^{2\beta}\), with comparison constants permitted to depend on that fixed noise level. Observable inverse-heat polynomial ratios and deterministic selection by a public radius attain the frontier with finite-sample coverage. Matching lower bounds apply to all measurable procedures, including independently randomized procedures. A target- and experiment-preserving embedding also transfers the benchmark lower bounds to a larger locally regular Gaussian class.
\end{abstract}

\section{Introduction}

This paper establishes a common minimax precision frontier for estimation and honest inference at a latent regression discontinuity cutoff. Treatment assignment records which side of the latent cutoff an individual occupies, while the researcher observes a score contaminated by Gaussian measurement error. On a specified bounded-support benchmark, we characterize both worst-case expected absolute estimation error and worst-case expected length of connected confidence intervals with uniform \(90\%\) coverage. The comparison holds simultaneously for every sample size \(n\ge2\), where \(n\) denotes the number of independent observations, and every known noise standard deviation \(\sigma\in[0,1]\), with constants independent of sample size and noise scale.

The experiment observes treatment, a binary outcome, and the contaminated score. Assignment follows the latent score, and the target is the difference between the two conditional potential-outcome means at the latent cutoff. Recorded assignment therefore separates the sample into two true-side arms, each of which presents an endpoint estimation problem. This structure connects continuity-based regression discontinuity \citep{Hahn2001,CattaneoTitiunik2022Survey} to statistical recovery under measurement error. In the closely related observed-treatment experiment, \citet[Assumptions 1Y, 2C, and 7; Proposition 4(b)]{PeiShen2017DevilTails} establish identification under Gaussian measurement error and local density positivity. Our question concerns the precision available from a finite sample under quantitative restrictions on the latent law.

Those restrictions define the domain of the matched frontier. The latent score has known support \([-1,1]\) and a density continuous on that interval, bounded everywhere between \(1/4\) and \(3/4\). Both potential outcomes are binary. Their specified continuous conditional means also lie between \(1/4\) and \(3/4\), and each has full-interval Hölder seminorm at most two for a known exponent \(\beta\in(0,1]\), the smoothness parameter. The standardized Gaussian error is independent of the latent score and both potential outcomes. The error distribution, its scale, the support, and the smoothness exponent are public inputs. These conditions specify the benchmark in \cref{obj:def:model}; the uniform estimation and coverage criteria range over every law in that class.

Our first result constructs an observable estimator and a finite-sample honest interval on this benchmark. Positive polynomial weights concentrate near the cutoff within each reflected treatment arm. A finite inverse-heat transform corrects these weights for the known Gaussian contamination. Outcome-marked and unmarked sample averages then estimate the corresponding latent weighted averages, and stabilized, projected ratios recover the arm means. A public radius combines a Hölder approximation bound with a Gaussian second-moment bound. Deterministic minimization of that radius over a finite candidate set selects the estimator and interval. \Cref{obj:thm:finite-certificate} establishes uniform coverage of at least \(0.9\), bounds worst expected absolute error by the selected radius, and bounds worst expected interval length by twice that radius.

Our second result shows that this construction attains the optimal order for both decision criteria. \Cref{obj:thm:uniform-frontier} characterizes the common rate through a unique scalar resolution equation, with comparison constants depending only on the fixed Hölder exponent. Matching lower bounds concern the complete observed experiment and apply to every measurable estimator and every admissible honest connected interval, including procedures with independent randomization. The witnesses combine full-interval Hölder regularity, separation of the cutoff contrast, and polynomial cancellation under Gaussian convolution.

The resolution equation distinguishes three regimes. When the noise scale is at most \(n^{-1/(2\beta+1)}\), both criteria have the direct-experiment order
\[
n^{-\beta/(2\beta+1)}.
\]
On the intermediate branch specified in \cref{obj:thm:uniform-frontier}, their common order is
\[
\frac{\sigma^\beta}
{[1+\log(n\sigma^{2\beta+1})]^{3\beta/2}}.
\]
For each fixed positive noise scale, the compact-support branch eventually gives
\[
\left\{\frac{\log\log(en)}{\log(en)}\right\}^{2\beta},
\]
with comparison constants in this fixed-noise specialization permitted to depend on that noise level. The finite-sample root characterization covers both interfaces, zero noise, and every shrinking-noise sequence. For example, with Lipschitz conditional means and \(\sigma=n^{-1/6}\), the intermediate branch applies for every \(n\ge2\).

The construction draws on compact-support deconvolution and Gaussian moment inversion. \citet[Theorem~1(a), fixed-support case]{Meister2007CompactSupport} analyzes density recovery under known compact support, while \citet[Theorem~1; Section~4.1]{WuYang2020DenoisedMoments} uses Gaussian polynomial inversion to estimate finite mixing distributions. Here polynomial inversion targets positive endpoint averages that are converted into conditional means through arm ratios. The noise-scale comparison is also related to shrinking-noise density estimation \citep[Theorems~2.1 and 2.2]{Delaigle2008ShrinkingNoise} and distribution estimation under Wasserstein loss \citep[Section~2, Theorems~2.4--2.6]{DurotMukherjee2026VanishingNoise}. Our target and losses give these comparisons their specific decision-theoretic meaning. The honest interval construction follows the smoothness-based principle of combining worst-case approximation bias with stochastic uncertainty \citep{ArmstrongKolesar2018OptimalInference,ArmstrongKolesar2020SimpleHonest}.

The benchmark has two further interpretations. First, \cref{obj:prop:published-class-converse-transfer} embeds it into a larger Gaussian class with local density positivity and continuous conditional means near the cutoff, preserving both the target and the randomized observed experiment. This embedding transfers the benchmark lower bounds to the larger class. Second, a prospective eligibility-effect design links administrative true-side assignment, a contaminated eligibility score, and a prespecified binary outcome. Under externally justified Gaussian calibration, a known bounded latent range, independent sampling, and the benchmark's global density and conditional-mean restrictions, the results supply a precision benchmark for that design.

Supporting details for the main results appear in the appendix containing \cref{obj:thm:direct-reduction,obj:thm:uniform-frontier-resolution}.

The paper proceeds through related work; the experiment, assumptions, and decision criteria; observable estimation and finite-sample inference; uniform minimax rates across noise scales; and interpretation and extensions. The appendices develop the polynomial construction, observed-experiment lower bounds, rate comparisons and verification scope, and proofs of the main results.
